% ====================================================================
% DATEI: content/10_klevels_full.tex
% Vollständige K-Level-Definitionen (K0–K12)
% Ontologie der Kontinua — Core 1.3
% ====================================================================

\section{Vollständige Hierarchie der Kontinua \texorpdfstring{\texorpdfstring{$K_0$}{K\_0}–$K_{12}$}{K0–K12}}
\label{sec:klevels-full}

Dieses Kapitel enthält die vollständige, rekonstruierte Beschreibung aller
Kontinuumsniveaus \(K_0\)–\(K_{12}\), die Material aus den ursprünglich
ausführlich verteilten Quellen in Core~2.x sowie den Läufen Physik, Chemie,
Biologie, Kognition und Soziales integriert.
Im Gegensatz zu Abschnitt~\ref{sec:model}, der einen kompakten Überblick bot,
erweitert das vorliegende Kapitel jedes Niveau zur vollständigen strukturellen
genauen Definition.

Core~1.3 fixiert die Hierarchie als \(K_0\)–\(K_{12}\), weil das öffentliche
Quellkorpus bereits explizite Ober-Level-Module, niveaubezogen Übergänge und
Vorhersage-/Prozesshüllen für \(K_{11}\) und \(K_{12}\) enthält.
Würde die Hierarchie so behandelt, als ende sie bei \(K_{10}\), wäre die
Monographie intern inkonsistent.

Jedes Kontinuumsniveau ist als Tupel definiert durch
\[
    K_x = \big(\Omega_x, A_x, P_x(t), J_x(t), \Theta_x,
           \partial\Omega_x, C_x, k_x(t)\big).
\]
Niveaus unterscheiden sich durch die Art ihrer Achsen, Potentiale, Schwellenwerte,
Randgeometrie, Zyklen und die Struktur des Embedding-Raums
\(M_x\), der für deren Existenz erforderlich ist.

% ====================================================================
\subsection{Überblick über die vertikale Hierarchie}
% ====================================================================

Die vertikale Hierarchie der Kontinua ist monoton:
\[
    K_0 \rightarrow K_1 \rightarrow K_2 \rightarrow \dots
    \rightarrow K_{10} \rightarrow K_{11} \rightarrow K_{12},
\]
wo jeder Übergang der Entstehung neuer Achsen,
Schwellenwerte, Flüsse und Zyklen entspricht, die nicht auf ein niedrigeres
Niveau reduzierbar sind.

Eine globale Übersicht ist in Tabelle~\ref{tab:klevels-summary} dargestellt.

\begin{table}[h!]
\centering
\begingroup
\footnotesize
\RaggedRight
\setlength{\tabcolsep}{4pt}
\begin{tabular}{@{}L{0.10\textwidth}L{0.25\textwidth}L{0.49\textwidth}@{}}
\hline
Level & Dominante Achsen & Charakteristische Struktur \\
\hline
\(K_0\) & Differenzachse & Strukturelles Substrat \\
\(K_1\) & Eindimensionale geometrische Achse & Klassisches Kontinuum \\
\(K_2\) & Physikalische Feldachsen & Phasen, BKT, Masse, Kohärenz \\
\(K_3\) & Chemische Achsen & RAF-Netzwerke, Katalyse \\
\(K_4\) & Membranachsen & Osmotische, Krümmungs- und Permeabilitäts-Schwellen \\
\(K_5\) & Erregbarkeitsachsen & Proto-Spikes, Ionenflüsse \\
\(K_6\) & Kognitive Achsen & Bindung, innere Modelle \\
\(K_7\) & Normative/soziale Achsen & Institutionen, Vertrauen \\
\(K_8\) & Zivilisatorische Achsen & Infrastruktur, kollektive Zyklen \\
\(K_9\) & Theoretische Achsen & Paradigmen, Ontologien \\
\(K_{10}\) & Meta-theoretische Achsen & Rekursion von Modell-von-Modell \\
\(K_{11}\) & Meta-evolutionäre formale Achsen & Operatoren, Meta-Logiken, evolvierende Modellräume \\
\(K_{12}\) & Globale Kohärenzachsen & Cross-framework semantische und strukturelle Kohärenz \\
\hline
\end{tabular}
\endgroup
\caption{Globale Übersicht der Kontinua \texorpdfstring{\(K_0\)}{K0}–\(K_{12}\).}
\label{tab:klevels-summary}
\end{table}

Jedes Niveau erfordert einen Einbettungsraum
\[
   M_0 \subset M_1 \subset \dots \subset M_{10} \subset M_{11} \subset M_{12},
\]
mit Achsen \(A(M_x)\), die ausreichend sind, um die Kontinua dieses Niveaus zu tragen.

% ====================================================================
\subsection{Niveau \texorpdfstring{\texorpdfstring{$K_0$}{K\_0}}{K0}: Strukturelles Substrat}
% ====================================================================

\paragraph{Formale Daten.}
\[
    K_0 = (S,\Delta,\mathcal{C}),
\]
mit:
\begin{itemize}
    \item \(S\): eine Menge unterscheidbarer Zustände;
    \item \(\Delta\): eine strukturelle Differenzfunktion;
    \item \(\mathcal{C}\): eine strukturelle Relation (oder Relationenfamilie),
          die Unterscheidbarkeit bewahrt.
\end{itemize}

Auf dieser Ebene existieren weder Zeit, Geometrie, Flüsse noch Potentiale.

\paragraph{Schwelle.}
Existenz-Schwelle:
\[
   \Theta_0 = \varepsilon > 0,
   \qquad
   s_1 \neq s_2 \Rightarrow \Delta(s_1,s_2) \ge \varepsilon.
\]

\paragraph{Rolle.}
Sie liefert das logische Substrat der Unterscheidbarkeit;
kein Kontinuum kann ohne dieses Niveau existieren.

% ====================================================================
\subsection{Niveau \texorpdfstring{\texorpdfstring{$K_1$}{K\_1}}{K1}: Eindimensionale Kontinua}
% ====================================================================

\paragraph{Formale Daten.}
\[
    K_1 = (\Omega_1, A_1, P_1, J_1, \Theta_1,
           \partial\Omega_1, C_1, k_1).
\]

\begin{itemize}
    \item \(A_1\): eine geometrische Achse (eindimensionales Kontinuum);
    \item \(\Omega_1\): klassischer Konfigurationsraum
          (z.\,B. \(C^0(T,H^1(X))\cap C^1(T,L^2(X))\));
    \item \(P_1\): klassisches Energie-Funktional;
    \item \(J_1\): klassische dynamische Flussfunktion;
    \item \(\Theta_1\): Stabilitäts- und kritische Schwellenwerte.
\end{itemize}

\paragraph{Beispiele.}
Klassische Felder in einer Raumdimension, gekoppelte Oszillatoren,
Skalar-Diffusionsmodelle.

\paragraph{Zyklen.}
Periodische Umlaufbahnen, stabile Grenzzyklen.

% ====================================================================
\subsection{Niveau \texorpdfstring{\texorpdfstring{$K_2$}{K\_2}}{K2}: Physikalische Kontinua}
% ====================================================================

\paragraph{Achsen.}
Multidimensionale Feldachsen:
\begin{itemize}
    \item Raumachsen;
    \item interne Feldachsen (z.\,B. Eichfreiheitsgrade);
    \item Ordnungsparameterachsen (Kohärenz, Magnetisierung usw.).
\end{itemize}

\paragraph{Potentiale und Flüsse.}
Physikalische Energie-Funktionale, Feldgleichungen, Renormierungsflüsse.

\paragraph{Schwellenwerte.}
\begin{itemize}
    \item \(\Theta_{\rm crit}\): Phasengrenzen;
    \item \(\Theta_{\rm BKT}\): Entwirrungs-Schwelle von Wirbelpaaren;
    \item \(\Theta_{\rm mass}\): Kohärenz-Schwelle für Massenerzeugung;
    \item \(\Theta_{\rm death}\): Konfinierung oder Dekopplung.
\end{itemize}

\paragraph{Zyklen.}
Topologische Solitonen, Wirbelpaare, feldtheoretische Zyklen.

% ====================================================================
\subsection{Niveaus \texorpdfstring{\texorpdfstring{$K_3$}{K\_3}–\texorpdfstring{$K_4$}{K\_4}}{K3–K4}: Chemische und Protokellulare Kontinua}
% ====================================================================

\subsubsection{Niveau \texorpdfstring{\(K_3\)}{K3}: Chemische Kontinua}

\paragraph{Achsen.}
Achsen chemischer Konzentrationen, Umgebungsachsen (pH, Salzgehalt, Temperatur).

\paragraph{Schwellenwerte.}
\begin{itemize}
    \item RAF-Existenz-Schwelle (Abschluss);
    \item katalytische Aktivierungsschwellen;
    \item Perkolationsschwelle für Reaktionsnetzwerke.
\end{itemize}

\paragraph{Zyklen.}
Chemische Zyklen, RAF-Abschluss-Zyklen.

\subsubsection{Niveau \texorpdfstring{\(K_4\)}{K4}: Protokellulare Kontinua}

\paragraph{Achsen.}
Membranachsen: Krümmung, Oberflächenspannung, Permeabilität, Ladung.

\paragraph{Schwellenwerte.}
\begin{itemize}
    \item osmotische Schwelle \(\Theta_{\rm grad}\);
    \item Permeabilitäts-Schwelle \(\Theta_{\rm perm}\);
    \item Membranruptur-Schwelle \(\Theta_{\rm mem}\);
    \item Krümmungsschwelle \(\Theta_{\rm curv}\).
\end{itemize}

\paragraph{Zyklen.}
Metabolische Zyklen, Membran-Wachstumszyklen, Gradienten-Erhaltung-Zyklen.

\paragraph{Randgeometrie.}
Detailliertes Patch-Modell:
Patch-Zustände \(\sigma_i\in\{L_\alpha,L_\beta,L_o,L_f,L_b\}\),
lokale Schwellenwert-Vektoren \(\Theta_i\),
Patch-Interaktionen und Rissdynamik.

% ====================================================================
\subsection{Niveau \texorpdfstring{\texorpdfstring{$K_5$}{K\_5}}{K5}: Frühneuronale und bioelektrische Kontinua}
% ====================================================================

\paragraph{Achsen.}
Erregbarkeitsachsen:
\begin{itemize}
    \item Membranpotentialachse \(\Delta V\);
    \item Ionenkonzentrationsachsen;
    \item Kanalzustandsachsen.
\end{itemize}

\paragraph{Potentiale und Flüsse.}
\begin{itemize}
    \item \(P_5\): elektrochemische Potentiale;
    \item \(J_5\): Ionenflüsse, Gate-Variablen-Flüsse.
\end{itemize}

\paragraph{Schwellenwerte.}
\begin{itemize}
    \item Erregbarkeitsschwelle \(\Theta_{\rm exc}\);
    \item Refraktär-Schwelle \(\Theta_{\rm refr}\);
    \item Gradienten-Schwelle \(\Theta_{\rm grad}\);
    \item Ausfalls-Schwelle \(\Theta_{\rm death}\) (Verlust der Erregbarkeit).
\end{itemize}

\paragraph{Zyklen.}
Proto-Spike-Zyklus \(C_{\rm spike}\),
frühe oszillatorische Zyklen,
Patch-Ebene-Erregungsschleifen.

% ====================================================================
\subsection{Niveau \texorpdfstring{\texorpdfstring{$K_6$}{K\_6}}{K6}: Kognitive Kontinua}
% ====================================================================

\paragraph{Achsen.}
\begin{itemize}
    \item Repräsentationale Achsen;
    \item Bindungsachsen;
    \item Konzeptuelle Achsen;
    \item Prädiktive Achsen.
\end{itemize}

\paragraph{Potentiale.}
\begin{itemize}
    \item semantische Energie-Potentiale;
    \item Vorhersagefehler-Potentiale;
    \item Gedächtnisstabilitäts-Potentiale.
\end{itemize}

\paragraph{Schwellenwerte.}
\begin{itemize}
    \item Bindungs-Schwelle \(\Theta_{\rm bind}\);
    \item Vorhersage-Schwelle \(\Theta_{\rm pred}\);
    \item Kohärenz-Schwelle \(\Theta_{\rm coh}\);
    \item expressive Schwelle \(\Theta_{\rm expr}\).
\end{itemize}

\paragraph{Zyklen.}
Kognitive Zyklen:
\begin{itemize}
    \item Aufmerksamkeitszyklus;
    \item Vorhersagezyklus;
    \item Gedächtniskonsolidierungszyklus;
    \item Modellaktualisierungszyklus.
\end{itemize}

% ====================================================================
\subsection{Niveaus \texorpdfstring{\texorpdfstring{$K_7$}{K\_7}–\texorpdfstring{$K_8$}{K\_8}}{K7–K8}: Soziale und zivilisatorische Kontinua}
% ====================================================================

\subsubsection{Niveau \texorpdfstring{\(K_7\)}{K7}: Soziale Kontinua}

\paragraph{Achsen.}
\begin{itemize}
    \item normative Achsen;
    \item Vertrauenskorridore;
    \item Rollenachsen;
    \item institutionelle Achsen.
\end{itemize}

\paragraph{Schwellenwerte.}
\begin{itemize}
    \item Vertrauens-Schwelle \(\Theta_{\rm trust}\);
    \item Legitimitäts-Schwelle \(\Theta_{\rm leg}\);
    \item Rollen-Kohärenz-Schwelle \(\Theta_{\rm role}\).
\end{itemize}

\paragraph{Zyklen.}
Normative Zyklen, institutionelle Zyklen, Koordinationszyklen.

\subsubsection{Niveau \texorpdfstring{\(K_8\)}{K8}: Zivilisatorische Kontinua}

\paragraph{Achsen.}
Makroachsen:
\begin{itemize}
    \item infrastrukturelle Achsen;
    \item Energieflussachsen;
    \item Kommunikationsachsen;
    \item Komplexitätsachsen.
\end{itemize}

\paragraph{Schwellenwerte.}
\begin{itemize}
    \item systemische Spannungs-Schwelle;
    \item Kapazitäts-Schwelle;
    \item Kohäsions-Schwelle.
\end{itemize}

\paragraph{Zyklen.}
Energiezyklen, Wirtschaftszyklen, institutionelle Megazyklen.

% ====================================================================
\subsection{Niveaus \texorpdfstring{\texorpdfstring{$K_9$}{K\_9}–$K_{12}$}{K9–K12}: Theoretische, meta-theoretische und obere Kohärenz-Kontinua}
% ====================================================================

\subsubsection{Niveau \texorpdfstring{\(K_9\)}{K9}: Theoretische Kontinua}

\paragraph{Achsen.}
\begin{itemize}
    \item konzeptuelle Achsen;
    \item formale Achsen;
    \item Modellierungsachsen.
\end{itemize}

\paragraph{Schwellenwerte.}
\begin{itemize}
    \item Konsistenz-Schwelle;
    \item Kohärenz-Schwelle;
    \item expressive Schwelle \(\Theta_{\rm expr}\).
\end{itemize}

\paragraph{Zyklen.}
Theorie-Evolutionszyklen, Paradigmenzyklen.

\subsubsection{Niveau \(K_{10}\): Meta-theoretische Kontinua}

\paragraph{Achsen.}
\begin{itemize}
    \item meta-linguistische Achsen;
    \item funktoriale Achsen;
    \item Modell-von-Modell-Achsen.
\end{itemize}

\paragraph{Schwellenwerte.}
\begin{itemize}
    \item Meta-Kohärenz-Schwelle;
    \item Selbstreferenz-Schwelle;
    \item Meta-Expressivitätsschwelle.
\end{itemize}

\paragraph{Zyklen.}
Funktoriale Zyklen, Rekursionzyklen der Meta-Theorie.

\subsubsection{Niveau \(K_{11}\): Rekursive Meta-theoretische Kontinua}

\paragraph{Achsen.}
\begin{itemize}
    \item Operatorfamilienachsen;
    \item meta-logische Achsen;
    \item Achsen zur Transformation von Modellräumen;
    \item Achsen rekursiver struktureller Reflexion.
\end{itemize}

\paragraph{Schwellenwerte.}
\begin{itemize}
    \item Operator-Kohärenz-Schwelle;
    \item rekursive Abschluss-Schwelle;
    \item Meta-Raum-Kompatibilitäts-Schwelle.
\end{itemize}

\paragraph{Rolle.}
\(K_{11}\) ist erforderlich, sobald Theorien, Operator-Suiten und Modellräume
selbst als sich entwickelnde strukturelle Objekte und nicht als statische
Behälter behandelt werden. Es ist daher die erste Ebene, auf der der Rahmen
Meta-Evolution explizit untersucht.

\paragraph{Zyklen.}
Operator-Update-Zyklen, Aktualisierungszyklen der Metalogik,
Modellraum-Rekursionszyklen.

\subsubsection{Niveau \(K_{12}\): Globale Kohärenz-Kontinua}

\paragraph{Achsen.}
\begin{itemize}
    \item semantische Kohärenzachsen;
    \item globale Kompatibilitätsachsen;
    \item achsenübergreifende Bindungsachsen.
\end{itemize}

\paragraph{Schwellenwerte.}
\begin{itemize}
    \item globale Konsistenz-Schwelle;
    \item semantische Kohärenz-Schwelle;
    \item obere Admissibilitäts-Schwelle.
\end{itemize}

\paragraph{Rolle.}
\(K_{12}\) ist die minimal definierte obere Schicht, wenn man globale
Kohärenzbedingungen in Familien von \(K_{11}\)-Strukturen diskutieren muss.
Ihre Rolle ist strukturell und begrenzend: Sie dokumentiert, dass die Hierarchie ein
explizites Oberproblem globaler Kohärenz besitzt, nicht bloß weiteres Wachstum der
gleichen unteren Dimensionen.

\paragraph{Zyklen.}
Globale Kohärenz-Wiederherstellungszyklen, semantische Ausrichtungszyklen,
Kreuz-Rahmen-Kooperationszyklen.

% ====================================================================
\subsection{Querebenen-Zusammenfassung}
% ====================================================================

Ein Kontinuum \(K_x\) existiert als lebendes Kontinuum genau dann, wenn
\[
   \Omega_x \neq \emptyset,
   \qquad
   A_x \subseteq A(M_x),
   \qquad
   \Theta_x \text{ ist innerhalb von } M_x \text{ erfüllbar.}
\]

Vertikale Kontinuitätsbedingungen:
\begin{itemize}
    \item \(A_x \subset A_{x+1}\);
    \item \(M_x \subset M_{x+1}\);
    \item \(\Theta_{x+1}\) verfeinert \(\Theta_x\);
    \item Flüsse \(J_x\) betten sich in \(J_{x+1}\) ein;
    \item Zyklen \(C_x\) bilden Teilstrukturen von \(C_{x+1}\).
\end{itemize}

Damit ist die rekonstruierte vollständige Hierarchie der Kontinua von
\(K_0\) bis \(K_{12}\) abgeschlossen.
